\documentclass{amsart}
\usepackage[T1]{fontenc}
\newtheorem{theorem}[subsection]{Theorem}

\usepackage{tikz}
\usepackage{hyperref}

\DeclareRobustCommand{\sagelogo}{%
  \begin{tikzpicture}[line width=.2ex,line cap=round,rounded corners=.01ex,baseline=-.2ex]
    \draw(0,0) -- (.75em,0) 
      -- (.75em,.7ex) -- (.25em,.7ex) 
      -- (.25em,.75\ht\strutbox) -- (.75em,.75\ht\strutbox);
    \draw(2em,0) -- (1.6em,0)
      -- (1.3em,.75\ht\strutbox) -- (.9em,.75\ht\strutbox)
      -- (.9em,0) -- (1.3em,0)
      -- (1.6em,.75\ht\strutbox) -- (2.1em,.75\ht\strutbox)
      -- (2.1em,-\dp\strutbox) -- (1.45em,-\dp\strutbox);
    \draw(3em,0) -- (2.25em,0)
      -- (2.25em,.75\ht\strutbox) -- (2.8em,.75\ht\strutbox)
      -- (2.8em,.7ex) -- (2.35em, .7ex);
  \end{tikzpicture}%
}


\title{On the classification of Rota-Baxter operators over $M_2(\mathbb{C})$}
\author{TODO}
\date{}

\begin{document}

\begin{abstract}
    A more refined classification of Rota-Baxter operators over $M_2(\mathbb{C})$ is presented.
\end{abstract}

\maketitle

\section{Raw results}
As proposed in \cite{TANG2014}, we utilized Computer Algebra systems \sagelogo \cite{sagemath}, {SINGULAR} \cite{singular} and {Macaulay2} \cite{Macaulay2} to classify Rota-Baxter operators. The supporting code is \href{https://github.com/AlexThePolygonal/Rota-Baxter-Cleaner-Code}{available on github}.

The projective variety of Rota-Baxter operators consists of 6 components: four curves $C_1, C_2, C_3, C_4$ and two surfaces $A$ and $B$.
\begin{itemize}
\item
Operators in C1 have the following shape:
$$\left(\begin{array}{rrrr}
-c & 0 & a & 0 \\
0 & -c & 0 & a \\
b & 0 & c & 0 \\
0 & b & 0 & c
\end{array}\right)$$
and their elements must lie in the ideal \begin{align*}
a b + c^{2} = 0 
\end{align*}
Its dimension as a projective variety is 1.
It is smooth.

\item
Operators in C2 have the following shape:
$$\left(\begin{array}{rrrr}
0 & 0 & -b & -a \\
b & a & 0 & 0 \\
0 & 0 & -a & -c \\
a & c & 0 & 0
\end{array}\right)$$
and their elements must lie in the ideal \begin{align*}
-a^{2} + b c = 0 
\end{align*}
Its dimension as a projective variety is 1.
It is smooth.

\item
Operators in C3 have the following shape:
$$\left(\begin{array}{rrrr}
-a & c & 0 & 0 \\
b & a & 0 & 0 \\
0 & 0 & -a & c \\
0 & 0 & b & a
\end{array}\right)$$
and their elements must lie in the ideal \begin{align*}
a^{2} + b c = 0 
\end{align*}
Its dimension as a projective variety is 1.
It is smooth.

\item
Operators in C4 have the following shape:
$$\left(\begin{array}{rrrr}
0 & -a & 0 & -c \\
0 & -c & 0 & -b \\
a & 0 & c & 0 \\
c & 0 & b & 0
\end{array}\right)$$
and their elements must lie in the ideal \begin{align*}
a b - c^{2} = 0 
\end{align*}
Its dimension as a projective variety is 1.
It is smooth.

\item
Operators in A have the following shape:
$$\left(\begin{array}{rrrr}
-h & -f & -b & -g \\
-b & -g & c & e \\
a & d & h & f \\
h & f & b & g
\end{array}\right)$$
and their elements must lie in the ideal \begin{align*}
-b f + g h = 0 \quad
-f^{2} + d g = 0 \quad
b d - f h = 0 \\
a g - f h = 0 \quad
a b - h^{2} = 0 \quad
a f - d h = 0 \\
e f + g^{2} = 0 \quad
b g + e h = 0 \quad
d e + f g = 0 \\
a e + b f = 0 \quad
-b e + c g = 0 \quad
c f + b g = 0 \\
b^{2} + c h = 0 \quad
c d + b f = 0 \quad
a c + b h = 0 
\end{align*}
Its dimension as a projective variety is 2.
It is smooth.

\item
Operators in B have the following shape:
$$\left(\begin{array}{rrrr}
-h & -f & -b & -g \\
a & h & c & b \\
-f & d & -g & e \\
h & f & b & g
\end{array}\right)$$
and their elements must lie in the ideal \begin{align*}
-b f + g h = 0 \quad
b e + g^{2} = 0 \quad
f g + e h = 0 \\
-e f + d g = 0 \quad
b d + f g = 0 \quad
f^{2} + d h = 0 \\
-b^{2} + c g = 0 \quad
c f - b h = 0 \quad
c e + b g = 0 \\
c d + b f = 0 \quad
a g - b h = 0 \quad
a b - c h = 0 \\
a f - h^{2} = 0 \quad
a e + b f = 0 \quad
a d + f h = 0 
\end{align*}
Its dimension as a projective variety is 2.
It is smooth.

\end{itemize}




All of the components are smooth. The curves $C_1, \ldots, C_4$ are all disjoint and do not intersect the surfaces. The surfaces intersect, and their intersection is a smooth curve $A \cap B$ operators on which have the following shape:
$$\left(\begin{array}{rrrr}
e & -d & -a & -e \\
-a & -e & b & a \\
-d & c & -e & d \\
-e & d & a & e
\end{array}\right)$$
and their elements must lie in the ideal \begin{align*}
a^{2} - b e = 0 \quad
a d + e^{2} = 0 \quad
a c + d e = 0 \\
d^{2} - c e = 0 \quad
b d + a e = 0 \quad
b c - e^{2} = 0 
\end{align*}


The transposition antihomomorphism exchanges the curves $C_1$ and $C_3$, $C_2$ and $C_4$, and the surfaces $A$ and $B$, leaving $A \cap B$ fixed.

The group $GL_n(\mathbb{C})$ acts naturally on the projective variety of Rota-Baxter operators. In a previous work by Benito, Gubarev and Pozhidaev \cite{Benito2018}, representatives of these orbits were found in the $2$-dimensional case. We compute the action of $SL_2(\mathbb{C})$ on these representatives and use it to describe the varieties further. Let the elements of $SL_2(\mathbb{C})$ be 
\[\begin{pmatrix}
  x & y\\
  z & t
\end{pmatrix}\]

\begin{itemize}
  \item The orbit of the points $M_1$ and $M_1^T$ both generate the singular locus $A \cap B$. The parametrization of the orbit by the $SL_2$-action is given by
  \begin{align*}
&  a = -x y^{3}\\
  &b = y^{4}\\
  &c = x^{4}\\
  &d = x^{3} y\\
  &e = x^{2} y^{2}
  &\end{align*}
  The stabilizer of the action is obvious.

  \item The orbit of the points $M_2$ and $M_2^T$ lie on the surfaces $A$ and $B$ respectively.
  The stabilizer of the action of $SL_2(\mathbb{C})$ on the point $M_2$ consists of all diagonal matrices in $SL_2(\mathbb{C})$
  The map giving the orbit for $A$ is given by 
  \begin{align*}
&  a = t x^{3}\\
  &b = t x y^{2}\\
  &c = -t y^{3}\\
  &d = -x^{3} z\\
  &e = y^{3} z\\
  &f = -x^{2} y z\\
  &g = -x y^{2} z\\
  &h = t x^{2} y
  &\end{align*}
  Note that since there is only one orbit in $A \setminus A\cap B$, then the mapping $SL_2(\mathbb{C}) \to A \setminus A \cap B$ is surjective and $A$ decomposes into two orbits.

 \item  The orbits of the points $M_3$ and $M_3^T$ are the curves $C_1$ and $C_3$ respectively. 
  The parametrization of the orbit given by the $SL_2(\mathbb{C})$-action on $C_1$ is given by
  \begin{align*}
  &a = -y^2\\
  &b = x^2\\
  &c = xy
  \end{align*}

  \item  The orbits of the points $M_4$ and $M_4^T$ are the curves $C_4$ and $C_2$ respectively. 
  The parametrization of the orbit given by the $SL_2(\mathbb{C})$-action on $C_4$ is given by
  \begin{align*}
  &a = -x^2\\
  &b = -y^2\\
  &c = -xy
  \end{align*}
\end{itemize}

Note that the action of $SL_2(\mathbb{C})$ on the orbits of the curves $C_1, \ldots, C_4, A\cap B$ gives their rational parametrisations, proving that they are isomorphic to $\mathbb{P}^1$.

The classes described by Tang et al. \cite{TANG2014} correspond to these as follows:
\begin{itemize}
    \item $C_2$: $R_1$,$R_{13}$
    \item $C_3$: $R_2$,$R_8$
    \item $B$: $R_3$,$R_{11}$,$R_{12}$,$R_{16}$,$R_{19}$,$R_{21}$
    \item $C_1$: $R_4$,$R_{15}$
    \item $C_4$: $R_5$,$R_9$
    \item $A$: $R_6$,$R_7$,$R_{10}$,$R_{14}$,$R_{17}$,$R_{18}$,$R_{20}$
\end{itemize}




\section{Limitations}


\bibliographystyle{plain}
\bibliography{annot.bib}

\end{document}